\paragraph{Explain operating cash from profit under explicit classifications.}
The preceding schedules leave a broader reporting question: how does each
recorded event reconcile profit with operating cash? Consider a separate
synthetic sale of plant for USD~12 when its carrying amount is USD~10. The
journal debits cash by 12, credits plant by 10 and credits gain income by 2.
With the receipt explicitly classified as investing, operating cash is zero
although profit is 2. An indirect explanation must remove that gain as well as
account for changes in receivables, inventory and unpaid operating obligations.
Balancing an unexplained difference would conceal which interpretation was used.

\texttt{DeclaredIndirectCash/2} supplies this explanation for a selected successful
report in a verified company program. Its policy identifies each used method
by its exact definition hash and classifies every noncash balance-sheet account
and every income or expense account that the method actually posts to. A name
such as \emph{disposal} cannot substitute for those classifications. Working-capital
movements, noncash profit adjustments, other nonoperating balance-sheet movements
and direct owner movements are separate balance-sheet categories. Income and
expense postings are separately designated operating or nonoperating. These are
supplied premises for the explanation; the report preserves the financial
statements and cash categories already produced by the company program.

To see why the explanation reconciles, first consider one event. All the
following amounts are exact monetary scalars in the book's currency, and signed
journal movements use the existing debit-positive convention. Let $P$ be profit,
the negative sum of the event's income and expense postings. Let $W$ be the
negative sum of postings classified as working-capital movements, and $N$ the
negative sum classified as noncash profit adjustments. Thus an increase in a
receivable contributes negatively to $W$, while an increase in an operating
payable contributes positively. Let $X$ be the signed sum of the remaining
noncash balance-sheet postings, those classified nonoperating or owner-related.
Finally, let $C_o$, $C_i$ and $C_f$ denote net operating, investing and financing
cash receipts; payments are negative. The balanced journal implies
\[
 C_o+C_i+C_f-P-W-N+X=0.
\]
This identity alone leaves the nonoperating part of profit inside $P$. Define
$R$ as the signed sum of income and expense postings explicitly classified as
nonoperating. A credited disposal gain therefore gives a negative $R$, whereas
a debited disposal loss gives a positive $R$. The policy must satisfy a second,
event-level check:
\[
 X+R+C_i+C_f=0.
\]
This checks that the excluded balance-sheet and profit movements account for
the investing and financing cash. Substituting
$X=-R-C_i-C_f$ in the balanced-journal identity gives
$C_o-P-W-N-R=0$, hence
\[
 C_o=P+W+N+R.
\]
The implementation checks both relations for every included event, then sums
the contributions over the selected period. Period profit and operating cash
must agree with the retained income and direct cash-flow statements. No
unexplained remainder becomes an adjustment.

In the plant-sale example, $P=2$, $W=N=0$, $X=-10$, $R=-2$ and $C_i=12$.
The exclusion check is $-10-2+12=0$, and operating cash is $2-2=0$.
If the same plant instead sells for USD~8, profit is $-2$, the loss adjustment
is $+2$, and $-10+2+8=0$ still reconciles the investing receipt. Depreciation
of USD~5 gives $P=-5$ and $N=5$, with no cash movement. An acquisition financed
directly by debt can likewise have offsetting nonoperating balance-sheet
postings and no cash receipt or payment; its gross postings remain visible.

Return to the combined operating and financing example. When interest payments
are operating, profit is USD~8.8 and the working-capital adjustment is
USD~$-5.8$, including the increase in unpaid interest. The bridge gives
USD~3 of operating cash. Under the separately supplied financing-interest
classification, the payable movements are nonoperating and the expense is
excluded from operating profit: $W=-8$ and $R=6.2$. The same profit then gives
$8.8-8+6.2=7$ dollars of operating cash. The USD~4 difference is the interest
payment moved between cash categories. Total cash and profit are unchanged;
neither reconciliation selects which policy is applicable.

The selected period begins at the original opening or the latest preceding
successful carry. A stopped program can explain a report already present in
its successful prefix, but cannot supply a later missing report. Opening,
closing, requests and complete results must reproduce exactly before the
explanation is accepted. The policy must cover exactly the methods used in
that period and their gross posting accounts, including offsetting postings
with zero net change. Missing or extra classifications refuse the explanation.
One account has one category within a method; mixed treatment of that account's
individual postings requires a further specification.

The finite checks cover both charts, period carry and a subsequent collection,
noncash acquisition, depreciation, gross offsetting postings, and disposal gains
and losses. They also preserve the earlier supported financial reports and
operating-cash totals. Algebraically consistent decompositions can nevertheless
disagree about whether a movement is working capital or another noncash
adjustment. Classification admission therefore remains false. The result
explains the selected recorded events under stated premises; it does not prove
that those premises satisfy an accounting standard, that every required event
was recorded, or that the full note package is complete.
